\section{Implementation and Reproducibility}
\label{app:protocol}
\subsection{Graph preprocessing and fixed checkpoints}
Table~\ref{tab:datasets} records the dataset dimensions from the execution plan. Features use the dataset loader's row normalization. Support pools and query sets come from the official training and test masks; validation labels are not used for fitting or selection. All other nodes remain visible as unlabeled graph context. Each task contains all classes of its graph, with the requested number of support nodes per class.
\begin{table}[!htbp]
\centering
\caption{\lang{Dataset dimensions and trainable prompt counts from the recorded protocol.}{记录协议中的数据维度及可学习提示参数量。}}
\label{tab:datasets}
\small
\setlength{\tabcolsep}{4pt}
\begin{tabular}{lrrrr}
\toprule
\lang{Dataset}{数据集} & \lang{Nodes}{节点} & \lang{Features}{特征} & \lang{Classes}{类别} & \lang{Prompt parameters}{提示参数} \\
\midrule
Cora & 2708 & 1433 & 7 & 896 \\
CiteSeer & 3327 & 3703 & 6 & 768 \\
PubMed & 19717 & 500 & 3 & 384 \\
\bottomrule
\end{tabular}
\end{table}


The encoder consists of two GCNConv layers of width 128 with LeakyReLU between them. The first layer includes the pretrained affine transformation and graph aggregation; the prompt is added to its output before the nonlinearity. The second convolution has no additional activation in the evaluated encoder. The same checkpoint is used by every comparison method on a dataset. Checkpoint, graph, and split hashes are preserved in the accompanying result records. The recorded pretraining family is \texttt{Edgepred\_Gprompt}; the pretraining random seed is unavailable.

\subsection{Training and prediction details}
The main text describes the active forward computation. For reproducibility, Adam uses $\beta_1=0.9$, $\beta_2=0.999$, $\epsilon=10^{-8}$, learning rate 0.01, and weight decay zero. The prompt regularizer is the mean of squared token entries multiplied by $10^{-3}$. It is separate from optimizer weight decay. The channel scale uses the sample standard deviation, with denominator $n-1$. Euclidean normalization uses the implementation's $10^{-12}$ norm floor. No distillation or query-label objective is active.

For the support loss, let $\bar z_i=\norm{Z_i}$ and $m_c=\sum_{j\in S_c}\bar z_j$. At $k>1$, the normalized prototype used for training support node $i$ is
\begin{equation}
p_c^{(-i)}=\norm{\frac{m_c-\mathbb I[y_i=c]\bar z_i}{k-\mathbb I[y_i=c]}},\qquad
\ell_{ic}=\bar z_i^\top p_c^{(-i)}/0.2.
\label{eq:loo}
\end{equation}
For $k=1$, the inclusive class prototype $\norm{m_c/k}$ replaces $p_c^{(-i)}$. During prediction, define
\begin{equation}
z_i^*=\norm{Z_i-\frac1n\sum_{j\in V}Z_j},\quad
p_c^*=\norm{\frac1k\sum_{i\in S_c}z_i^*},\quad
\widehat y_i=\arg\max_c (z_i^*)^\top p_c^*.
\label{eq:predict}
\end{equation}
The probability scores for all five metrics are obtained by applying softmax to cosine logits divided by 0.2. Training does not apply the output centering used at prediction. The readout is a fixed implementation choice, not an experimental factor in this paper.

\paragraph{Procedure.}
Given a graph, a fixed checkpoint, and support labels:
\begin{enumerate}
\item Normalize adjacency and compute Eq.~\eqref{eq:attributes} and the first-layer preactivation.
\item Initialize support prototypes with Eq.~\eqref{eq:anchor} and apply Eqs.~\eqref{eq:assignment}--\eqref{eq:refinement} three times.
\item Compute the fixed route with Eq.~\eqref{eq:routing}; for the shared-prompt ablation, replace it by $B_{ic}=1/C$.
\item Initialize $P=0$ and cache the hidden-channel scales. Keep the encoder and route fixed.
\item For 100 steps, run the prompted encoder, form support logits, and update $P$ using Eq.~\eqref{eq:loss}.
\item Compute final prototype scores using Eq.~\eqref{eq:predict}; seal predictions before scoring with query labels.
\end{enumerate}

\subsection{Baseline fitting and seed control}
The support-draw seeds are 30093000 through 30093019 and training seeds are 0, 1, and 2. A draw uses a seeded permutation in each class's support pool, taking its first $k$ entries. All methods receive the same support identities for each draw. The no-training Frozen method retains the same three seed slots for alignment, but those slots do not represent three independent fits.

The learnable comparison methods use support cross-entropy. Full FT and the input-prompt methods include a task readout, whereas the GraphPrompt adaptation uses its prompted representation with prototype prediction. GPF adds a shared feature vector; GPF-plus learns 20 feature prompts and an input-dependent mixture. Their final one-shot parameters are learning rate 0.01 and weight decay 0.0001. For 3 and 5 shots, each of three support-only folds holds out one node per class. Four candidate configurations are ranked by mean held-out support cross-entropy, then the selected method is refit on the complete support set. No external validation-label budget is used. This protocol controls labeled information and final step count, but not total selection compute or readout parameterization.

\subsection{Condition-specific statistics}
For a fixed dataset, shot count, checkpoint, and method, let $a_{st}$ denote a metric with support seed $s$ and training seed $t$. We compute
\begin{equation}
\bar a_s=\frac13\sum_{t=1}^{3}a_{st},\quad
\widehat a=\frac1{20}\sum_{s=1}^{20}\bar a_s,\quad
\widehat\sigma=\sqrt{\frac1{19}\sum_{s=1}^{20}(\bar a_s-\widehat a)^2}.
\label{eq:stats}
\end{equation}
Tables report $100\widehat a\pm100\widehat\sigma$, not the standard error. For paired comparisons, $d_s=\bar a_s^{\mathrm{PRPrompt}}-\bar a_s^{\mathrm{reference}}$ is resampled as a support-draw block. The percentile bounds of 10,000 resampled means form the reported 95\% interval. Existing intervals are reused without modification. No metric average or mean effect is computed across different datasets, shot counts, or methods.

Accuracy, Macro-F1, Micro-F1, macro one-vs-rest ROC-AUC, and macro average precision are retained in the data package. In this single-label multiclass setting, Micro-F1 coincides with Accuracy; it is included for completeness. The fixed query set is shared across support draws, so intervals concern support selection and training variation conditional on that graph, not independent samples of graphs.

\subsection{Environment, provenance, and numerical qualification}
The execution plan records Python 3.11.16, PyTorch 2.11.0+cu130, and CUDA 13.0 in the reused ProG environment. Per-unit records identify CUDA devices and fitting durations. An isolated total GPU-hour budget for the present reporting scope and a verified accelerator inventory are not provided by the extracted package; no efficiency comparison is inferred from these results.

The tables derive from the sealed three-dataset package. It contains 3,240 main-comparison units and 540 additional shared-prompt units; \method{} contributes 540 units. Source-metric checksums and paired graph/checkpoint/support/query identities were verified during extraction. The source numerical audits report exact readout reproduction from saved representations for all 3,780 selected units, passing gate/forward tolerances, and no changed query classes in fresh forwards. These are archived verification results, not newly rerun experiments, and passing a tolerance does not imply bitwise-identical CUDA arithmetic.

Source code and the table data are available in the accompanying repository at \url{https://github.com/Mes-sala/PRPrompt}. The source revision of the recorded experiments is \texttt{0328a2bb68\allowbreak f249a34c89\allowbreak 557178b623\allowbreak 3daf5ce20d}. The paper package includes table-generation and static-validation scripts, hashes, and bibliographic provenance. Model checkpoints and optimizer states remain on the experiment server and are excluded from the local paper package. This manuscript reports existing experiments; no new training or bootstrap analysis was performed for its preparation.

\clearpage
\section{Additional Condition-Wise Metrics}
\label{app:more}
The following tables extend the main comparison with the remaining stored metrics under the same seed aggregation rule. Every row is a separate dataset/shot condition.
\begin{table}[htbp]
\centering
\caption{\lang{Micro-F1 (\%, higher is better). Main comparison. Within each dataset/shot/method, training seeds are averaged within each support draw; mean and sample standard deviation are then computed over 20 support draws. Bold marks the largest mean in that row. $\dagger$: common-protocol adaptation.}{Micro-F1（\%，越高越好）。主比较。固定数据集、shot 与方法，在每个支持抽样内先平均训练 seed，再对 20 个支持抽样计算均值与样本标准差。粗体为该行最大均值。 $\dagger$：统一协议适配实现。}}
\label{tab:main-micro_f1}
\small
\setlength{\tabcolsep}{4pt}
\resizebox{\linewidth}{!}{%
\begin{tabular}{lrrrrrr}
\toprule
\lang{Dataset / shot}{数据集 / shot} & \frozenname & \ftname & GPF & GPF-plus & \gpromptname & \method \\
\midrule
Cora / 1 & \meanstd{23.17}{5.48} & \meanstd{45.41}{10.18} & \meanstd{15.20}{3.24} & \meanstd{30.80}{6.64} & \meanstd{21.88}{4.62} & \bestmeanstd{54.59}{9.18} \\
Cora / 3 & \meanstd{29.42}{3.47} & \meanstd{61.54}{6.78} & \meanstd{24.32}{4.66} & \meanstd{40.21}{4.80} & \meanstd{27.75}{3.64} & \bestmeanstd{70.76}{2.76} \\
Cora / 5 & \meanstd{33.56}{5.03} & \meanstd{67.50}{5.06} & \meanstd{28.17}{5.54} & \meanstd{53.14}{5.55} & \meanstd{31.36}{3.61} & \bestmeanstd{73.02}{1.95} \\
CiteSeer / 1 & \meanstd{21.14}{2.87} & \meanstd{35.76}{6.80} & \meanstd{17.10}{2.64} & \meanstd{21.40}{4.60} & \meanstd{19.53}{3.34} & \bestmeanstd{47.12}{6.31} \\
CiteSeer / 3 & \meanstd{25.29}{2.16} & \meanstd{42.22}{5.88} & \meanstd{20.36}{2.02} & \meanstd{28.92}{2.68} & \meanstd{24.83}{2.77} & \bestmeanstd{61.20}{4.98} \\
CiteSeer / 5 & \meanstd{27.38}{2.84} & \meanstd{47.89}{4.85} & \meanstd{21.40}{2.57} & \meanstd{31.82}{2.60} & \meanstd{27.66}{2.49} & \bestmeanstd{66.17}{2.16} \\
PubMed / 1 & \meanstd{55.44}{11.42} & \meanstd{48.77}{8.54} & \meanstd{36.29}{5.79} & \meanstd{41.12}{6.83} & \meanstd{50.24}{12.95} & \bestmeanstd{58.70}{9.78} \\
PubMed / 3 & \meanstd{65.55}{8.11} & \meanstd{59.75}{7.97} & \meanstd{51.28}{6.93} & \meanstd{57.46}{4.96} & \meanstd{62.47}{5.54} & \bestmeanstd{68.66}{4.58} \\
PubMed / 5 & \meanstd{68.37}{3.41} & \meanstd{63.96}{5.18} & \meanstd{55.83}{4.93} & \meanstd{61.37}{5.32} & \meanstd{66.94}{4.09} & \bestmeanstd{70.98}{2.46} \\
\bottomrule
\end{tabular}
}
\end{table}

\begin{table}[!htbp]
\centering
\caption{\lang{Macro ROC-AUC (\%, higher is better). Main comparison; seed aggregation follows Appendix~\ref{app:protocol}.}{Macro ROC-AUC（\%，越高越好）。主比较；seed 统计沿用附录~\ref{app:protocol}。}}
\label{tab:main-roc_auc_macro}
\small
\setlength{\tabcolsep}{4pt}
\resizebox{\linewidth}{!}{%
\begin{tabular}{lrrrrrr}
\toprule
\lang{Dataset / shot}{数据集 / shot} & \frozenname & \ftname & GPF & GPF-plus & \gpromptname & \method \\
\midrule
Cora / 1 & \meanstd{62.95}{4.93} & \meanstd{81.83}{5.73} & \meanstd{53.20}{2.50} & \meanstd{74.73}{4.00} & \meanstd{61.93}{4.85} & \bestmeanstd{83.06}{5.42} \\
Cora / 3 & \meanstd{68.95}{2.99} & \meanstd{89.94}{2.35} & \meanstd{64.61}{4.25} & \meanstd{78.13}{2.49} & \meanstd{68.05}{3.01} & \bestmeanstd{90.84}{1.97} \\
Cora / 5 & \meanstd{72.36}{2.85} & \bestmeanstd{92.58}{1.26} & \meanstd{68.65}{4.34} & \meanstd{85.66}{2.71} & \meanstd{71.01}{2.66} & \meanstd{92.47}{1.11} \\
CiteSeer / 1 & \meanstd{56.11}{3.24} & \meanstd{70.97}{4.63} & \meanstd{53.08}{1.84} & \meanstd{64.19}{4.07} & \meanstd{54.39}{4.21} & \bestmeanstd{74.71}{5.81} \\
CiteSeer / 3 & \meanstd{60.17}{2.06} & \meanstd{76.04}{3.00} & \meanstd{55.85}{1.91} & \meanstd{65.37}{3.23} & \meanstd{59.99}{2.60} & \bestmeanstd{82.82}{3.08} \\
CiteSeer / 5 & \meanstd{62.69}{2.35} & \meanstd{79.54}{2.22} & \meanstd{57.93}{1.77} & \meanstd{68.91}{2.38} & \meanstd{62.67}{2.33} & \bestmeanstd{85.86}{1.52} \\
PubMed / 1 & \meanstd{73.87}{11.03} & \meanstd{70.73}{9.76} & \meanstd{60.13}{8.31} & \meanstd{68.63}{6.61} & \meanstd{70.43}{11.06} & \bestmeanstd{76.48}{10.43} \\
PubMed / 3 & \meanstd{82.17}{6.59} & \meanstd{81.00}{3.74} & \meanstd{72.19}{4.91} & \meanstd{76.82}{4.24} & \meanstd{81.13}{4.42} & \bestmeanstd{86.29}{3.10} \\
PubMed / 5 & \meanstd{84.83}{2.59} & \meanstd{84.31}{2.45} & \meanstd{74.64}{4.10} & \meanstd{79.98}{2.95} & \meanstd{83.57}{2.81} & \bestmeanstd{87.03}{1.87} \\
\bottomrule
\end{tabular}
}
\end{table}

\begin{table}[!htbp]
\centering
\caption{\lang{Macro AP (\%, higher is better). Main comparison; seed aggregation follows Appendix~\ref{app:protocol}.}{Macro AP（\%，越高越好）。主比较；seed 统计沿用附录~\ref{app:protocol}。}}
\label{tab:main-average_precision_macro}
\small
\setlength{\tabcolsep}{4pt}
\resizebox{\linewidth}{!}{%
\begin{tabular}{lrrrrrr}
\toprule
\lang{Dataset / shot}{数据集 / shot} & \frozenname & \ftname & GPF & GPF-plus & \gpromptname & \method \\
\midrule
Cora / 1 & \meanstd{23.64}{4.26} & \meanstd{54.01}{10.10} & \meanstd{16.64}{1.47} & \meanstd{40.37}{5.88} & \meanstd{23.15}{3.89} & \bestmeanstd{59.04}{9.07} \\
Cora / 3 & \meanstd{29.20}{3.71} & \meanstd{70.20}{4.98} & \meanstd{24.89}{3.11} & \meanstd{44.43}{4.09} & \meanstd{27.06}{2.89} & \bestmeanstd{74.33}{3.38} \\
Cora / 5 & \meanstd{33.16}{4.26} & \meanstd{75.99}{3.16} & \meanstd{28.68}{3.71} & \meanstd{58.84}{6.14} & \meanstd{29.92}{3.11} & \bestmeanstd{77.64}{1.87} \\
CiteSeer / 1 & \meanstd{21.53}{2.32} & \meanstd{38.74}{5.84} & \meanstd{19.66}{1.43} & \meanstd{29.38}{4.20} & \meanstd{20.61}{2.27} & \bestmeanstd{45.99}{6.79} \\
CiteSeer / 3 & \meanstd{24.17}{1.59} & \meanstd{45.56}{4.79} & \meanstd{21.40}{1.15} & \meanstd{30.59}{3.23} & \meanstd{23.97}{1.89} & \bestmeanstd{59.15}{4.80} \\
CiteSeer / 5 & \meanstd{26.46}{2.07} & \meanstd{51.03}{4.40} & \meanstd{22.88}{1.45} & \meanstd{34.40}{2.70} & \meanstd{26.04}{1.65} & \bestmeanstd{64.02}{2.66} \\
PubMed / 1 & \meanstd{57.52}{10.93} & \meanstd{55.36}{10.55} & \meanstd{43.93}{6.89} & \meanstd{52.61}{8.08} & \meanstd{52.39}{9.87} & \bestmeanstd{61.58}{11.39} \\
PubMed / 3 & \meanstd{67.27}{8.59} & \meanstd{67.52}{6.82} & \meanstd{55.29}{5.93} & \meanstd{60.88}{5.58} & \meanstd{65.03}{5.79} & \bestmeanstd{74.22}{5.19} \\
PubMed / 5 & \meanstd{70.34}{4.13} & \meanstd{72.33}{5.00} & \meanstd{58.42}{5.12} & \meanstd{65.42}{4.92} & \meanstd{68.42}{3.67} & \bestmeanstd{75.35}{3.34} \\
\bottomrule
\end{tabular}
}
\end{table}

\clearpage
\section{Paired Differences and Uncertainty}
\label{app:paired}
Each cell reports \method{} minus the named reference in percentage points, with its unadjusted 95\% paired support-bootstrap interval. All entries preserve the original paired statistics. Tables group conditions for readability but do not average them.
\begin{table}[!htbp]
\centering
\caption{\lang{\method{} minus \frozenname: paired differences and 95\% intervals (percentage points).}{\method{} 相对 \frozenname 的配对差值与 95\% 区间（百分点）。}}
\label{tab:paired-FrozenPrototype}
\small
\setlength{\tabcolsep}{4pt}
\begin{tabular}{lrr}
\toprule
\lang{Dataset / shot}{数据集 / shot} & Accuracy & Macro-F1 \\
\midrule
Cora / 1 & $+31.42\;[+28.20,\,+34.59]$ & $+30.84\;[+27.77,\,+33.84]$ \\
Cora / 3 & $+41.34\;[+39.63,\,+43.04]$ & $+40.32\;[+38.64,\,+42.08]$ \\
Cora / 5 & $+39.46\;[+37.23,\,+41.61]$ & $+38.97\;[+36.88,\,+40.99]$ \\
CiteSeer / 1 & $+25.98\;[+23.72,\,+28.21]$ & $+23.87\;[+21.80,\,+25.77]$ \\
CiteSeer / 3 & $+35.91\;[+33.05,\,+38.39]$ & $+33.87\;[+31.34,\,+36.14]$ \\
CiteSeer / 5 & $+38.78\;[+37.37,\,+40.18]$ & $+36.58\;[+35.22,\,+37.95]$ \\
PubMed / 1 & $+3.26\;[+0.40,\,+6.22]$ & $+3.93\;[+0.86,\,+7.00]$ \\
PubMed / 3 & $+3.11\;[+0.38,\,+6.27]$ & $+5.02\;[+2.27,\,+8.44]$ \\
PubMed / 5 & $+2.61\;[+1.13,\,+4.22]$ & $+4.12\;[+2.53,\,+5.91]$ \\
\bottomrule
\end{tabular}
\end{table}

\begin{table}[!htbp]
\centering
\caption{\lang{\method{} minus \ftname: paired differences and 95\% intervals (percentage points).}{\method{} 相对 \ftname 的配对差值与 95\% 区间（百分点）。}}
\label{tab:paired-FullFineTune}
\small
\setlength{\tabcolsep}{4pt}
\begin{tabular}{lrr}
\toprule
\lang{Dataset / shot}{数据集 / shot} & Accuracy & Macro-F1 \\
\midrule
Cora / 1 & $+9.17\;[+6.58,\,+11.87]$ & $+8.97\;[+6.40,\,+11.77]$ \\
Cora / 3 & $+9.22\;[+6.75,\,+11.79]$ & $+8.44\;[+6.12,\,+10.90]$ \\
Cora / 5 & $+5.52\;[+3.73,\,+7.50]$ & $+5.05\;[+3.34,\,+6.86]$ \\
CiteSeer / 1 & $+11.36\;[+9.57,\,+13.18]$ & $+13.27\;[+11.17,\,+15.31]$ \\
CiteSeer / 3 & $+18.98\;[+16.78,\,+20.94]$ & $+19.76\;[+17.45,\,+22.00]$ \\
CiteSeer / 5 & $+18.28\;[+16.22,\,+20.13]$ & $+18.42\;[+16.07,\,+20.60]$ \\
PubMed / 1 & $+9.93\;[+6.77,\,+13.14]$ & $+12.86\;[+9.39,\,+16.37]$ \\
PubMed / 3 & $+8.90\;[+5.89,\,+12.26]$ & $+10.57\;[+7.08,\,+14.41]$ \\
PubMed / 5 & $+7.02\;[+4.96,\,+9.14]$ & $+7.86\;[+5.54,\,+10.28]$ \\
\bottomrule
\end{tabular}
\end{table}

\clearpage
\begin{table}[!htbp]
\centering
\caption{\lang{\method{} minus GPF: paired differences and 95\% intervals (percentage points).}{\method{} 相对 GPF 的配对差值与 95\% 区间（百分点）。}}
\label{tab:paired-GPF}
\small
\setlength{\tabcolsep}{4pt}
\begin{tabular}{lrr}
\toprule
\lang{Dataset / shot}{数据集 / shot} & Accuracy & Macro-F1 \\
\midrule
Cora / 1 & $+39.38\;[+35.60,\,+42.98]$ & $+46.27\;[+42.36,\,+49.82]$ \\
Cora / 3 & $+46.44\;[+44.21,\,+48.74]$ & $+48.25\;[+45.22,\,+51.53]$ \\
Cora / 5 & $+44.86\;[+42.52,\,+47.57]$ & $+46.17\;[+43.71,\,+49.32]$ \\
CiteSeer / 1 & $+30.02\;[+27.04,\,+33.00]$ & $+37.48\;[+34.93,\,+39.92]$ \\
CiteSeer / 3 & $+40.83\;[+38.11,\,+43.06]$ & $+44.38\;[+40.82,\,+47.90]$ \\
CiteSeer / 5 & $+44.77\;[+43.32,\,+46.16]$ & $+47.77\;[+44.95,\,+50.72]$ \\
PubMed / 1 & $+22.40\;[+18.82,\,+26.24]$ & $+31.76\;[+28.23,\,+35.27]$ \\
PubMed / 3 & $+17.37\;[+14.16,\,+21.02]$ & $+21.33\;[+17.31,\,+25.89]$ \\
PubMed / 5 & $+15.15\;[+13.30,\,+17.01]$ & $+17.97\;[+15.46,\,+20.67]$ \\
\bottomrule
\end{tabular}
\end{table}

\begin{table}[!htbp]
\centering
\caption{\lang{\method{} minus GPF-plus: paired differences and 95\% intervals (percentage points).}{\method{} 相对 GPF-plus 的配对差值与 95\% 区间（百分点）。}}
\label{tab:paired-GPFPlus}
\small
\setlength{\tabcolsep}{4pt}
\begin{tabular}{lrr}
\toprule
\lang{Dataset / shot}{数据集 / shot} & Accuracy & Macro-F1 \\
\midrule
Cora / 1 & $+23.78\;[+19.96,\,+27.47]$ & $+26.26\;[+22.23,\,+30.20]$ \\
Cora / 3 & $+30.55\;[+28.29,\,+32.52]$ & $+30.07\;[+28.04,\,+31.97]$ \\
Cora / 5 & $+19.88\;[+17.62,\,+22.16]$ & $+20.22\;[+18.02,\,+22.39]$ \\
CiteSeer / 1 & $+25.71\;[+22.89,\,+28.45]$ & $+31.43\;[+28.61,\,+34.12]$ \\
CiteSeer / 3 & $+32.27\;[+30.45,\,+34.03]$ & $+32.03\;[+30.10,\,+33.81]$ \\
CiteSeer / 5 & $+34.35\;[+33.05,\,+35.65]$ & $+33.57\;[+32.22,\,+34.87]$ \\
PubMed / 1 & $+17.57\;[+13.28,\,+21.86]$ & $+22.48\;[+18.10,\,+26.98]$ \\
PubMed / 3 & $+11.20\;[+9.47,\,+13.11]$ & $+13.74\;[+12.07,\,+15.56]$ \\
PubMed / 5 & $+9.61\;[+7.84,\,+11.49]$ & $+11.61\;[+9.80,\,+13.53]$ \\
\bottomrule
\end{tabular}
\end{table}

\clearpage
\begin{table}[!htbp]
\centering
\caption{\lang{\method{} minus \gpromptname: paired differences and 95\% intervals (percentage points).}{\method{} 相对 \gpromptname 的配对差值与 95\% 区间（百分点）。}}
\label{tab:paired-GraphPromptDiag}
\small
\setlength{\tabcolsep}{4pt}
\begin{tabular}{lrr}
\toprule
\lang{Dataset / shot}{数据集 / shot} & Accuracy & Macro-F1 \\
\midrule
Cora / 1 & $+32.70\;[+29.19,\,+36.29]$ & $+32.32\;[+28.93,\,+35.68]$ \\
Cora / 3 & $+43.01\;[+41.38,\,+44.73]$ & $+42.41\;[+40.61,\,+44.20]$ \\
Cora / 5 & $+41.66\;[+40.03,\,+43.20]$ & $+41.22\;[+39.64,\,+42.75]$ \\
CiteSeer / 1 & $+27.59\;[+24.97,\,+30.09]$ & $+26.26\;[+23.90,\,+28.62]$ \\
CiteSeer / 3 & $+36.37\;[+33.26,\,+38.82]$ & $+34.48\;[+31.80,\,+36.73]$ \\
CiteSeer / 5 & $+38.51\;[+37.14,\,+39.90]$ & $+36.31\;[+34.99,\,+37.65]$ \\
PubMed / 1 & $+8.45\;[+3.60,\,+13.47]$ & $+12.02\;[+7.33,\,+16.52]$ \\
PubMed / 3 & $+6.19\;[+4.04,\,+8.30]$ & $+8.27\;[+6.21,\,+10.36]$ \\
PubMed / 5 & $+4.03\;[+1.94,\,+6.36]$ & $+5.60\;[+3.49,\,+7.88]$ \\
\bottomrule
\end{tabular}
\end{table}

\begin{table}[!htbp]
\centering
\caption{\lang{\method{} minus \sharedname: paired differences and 95\% intervals (percentage points).}{\method{} 相对 \sharedname 的配对差值与 95\% 区间（百分点）。}}
\label{tab:paired-NoRouting}
\small
\setlength{\tabcolsep}{4pt}
\begin{tabular}{lrr}
\toprule
\lang{Dataset / shot}{数据集 / shot} & Accuracy & Macro-F1 \\
\midrule
Cora / 1 & $+31.96\;[+28.59,\,+35.21]$ & $+31.70\;[+28.29,\,+35.00]$ \\
Cora / 3 & $+39.06\;[+36.49,\,+41.44]$ & $+39.13\;[+36.60,\,+41.48]$ \\
Cora / 5 & $+36.94\;[+34.87,\,+39.06]$ & $+37.15\;[+35.09,\,+39.23]$ \\
CiteSeer / 1 & $+26.40\;[+24.07,\,+28.73]$ & $+24.69\;[+22.70,\,+26.65]$ \\
CiteSeer / 3 & $+37.18\;[+34.90,\,+39.33]$ & $+35.52\;[+33.31,\,+37.55]$ \\
CiteSeer / 5 & $+40.45\;[+39.06,\,+41.84]$ & $+38.48\;[+37.22,\,+39.72]$ \\
PubMed / 1 & $+4.85\;[+1.26,\,+8.26]$ & $+5.44\;[+1.84,\,+9.07]$ \\
PubMed / 3 & $+5.69\;[+3.25,\,+8.27]$ & $+7.27\;[+4.64,\,+10.29]$ \\
PubMed / 5 & $+3.02\;[+1.39,\,+4.51]$ & $+4.10\;[+2.56,\,+5.58]$ \\
\bottomrule
\end{tabular}
\end{table}

